Tall reinforced-concrete buildings in seismically active regions of the Himalayan foothills must
resist long-duration ground motions whose frequency content can excite higher vibration modes.
Code-based design relies on equivalent static forces and response spectra calibrated mostly on
shorter buildings. This thesis evaluates how well those procedures predict the demands on tall
core-wall buildings. It develops nonlinear finite-element models of three buildings of 20, 30
and 40 storeys, subjects them to a suite of 22 recorded and simulated ground motions, and compares
the resulting storey drifts and wall shear forces with code estimates. The response-spectrum
method underestimates the shear force at the base of the core wall by 25\% to 55\%, with the
error growing with height because of higher-mode contributions. A modified modal combination is
proposed that brings the estimate within 12\% of the nonlinear results.

\medskip\noindent\textbf{Keywords:} tall buildings, reinforced-concrete core walls, nonlinear
dynamic analysis, higher modes, response spectrum
